\documentclass[a4paper,12pt]{ltjsarticle}
\usepackage{luatexja}
\usepackage{amsmath,amssymb,amsthm}
\usepackage{mathtools}
\usepackage{tikz-cd}
\usepackage{tikz}
\usetikzlibrary{positioning,arrows.meta,calc,decorations.pathmorphing}
\usepackage{tcolorbox}
\tcbuselibrary{theorems,skins,breakable}
\usepackage{hyperref}
\usepackage{enumitem}
\usepackage{xcolor}

% 定理環境の設定
\theoremstyle{definition}
\newtheorem{definition}{定義}[section]
\newtheorem{theorem}[definition]{定理}
\newtheorem{lemma}[definition]{補題}
\newtheorem{proposition}[definition]{命題}
\newtheorem{corollary}[definition]{系}
\newtheorem{example}[definition]{例}
\newtheorem{remark}[definition]{注意}

% tcolorboxスタイル
\newtcolorbox{maintheorem}[1][]{
  colback=blue!5!white,
  colframe=blue!75!black,
  fonttitle=\bfseries,
  title=#1,
  breakable
}

\newtcolorbox{keyidea}[1][]{
  colback=green!5!white,
  colframe=green!75!black,
  fonttitle=\bfseries,
  title=#1,
  breakable
}

\newtcolorbox{computation}[1][]{
  colback=orange!5!white,
  colframe=orange!75!black,
  fonttitle=\bfseries,
  title=#1,
  breakable
}

% タイトル設定
\title{計算可能マルチンゲール理論と\\Optional Stopping Theorem}
\author{%
  Lean 4 形式化：CODEX (OpenAI) 生成コード\\[0.5em]
  \LaTeX 解説：Claude Code 生成ドキュメント%
}
\date{\today}

\begin{document}

\maketitle

\begin{abstract}
本稿は、離散有限標本空間上で完全に計算可能なマルチンゲール理論と、
その核心定理である Optional Stopping Theorem (OST) の形式化について解説する。
Lean 4 による実装は CODEX (OpenAI) により生成され、本 \LaTeX 解説文書は Claude Code により生成された。
有限性という制約の下で、確率論の本質的な構造を構造塔（Structure Tower）理論との対応において明らかにする。
\end{abstract}

\tableofcontents

\section{導入}

\subsection{背景と動機}

マルチンゲール理論は現代確率論の中核をなす分野であり、
金融数学、統計力学、情報理論など幅広い応用を持つ。
その中心的な定理の一つが \textbf{Optional Stopping Theorem (OST)} である：

\begin{maintheorem}[Optional Stopping Theorem（古典版）]
$(M_n)_{n \geq 0}$ をマルチンゲールとし、$\tau$ を有界停止時間とする。
このとき
\[
  \mathbb{E}[M_\tau] = \mathbb{E}[M_0].
\]
\end{maintheorem}

本稿では、この定理を\textbf{離散有限標本空間}という制限された設定で完全に形式化する。
この制約により：
\begin{itemize}
  \item すべての演算が\textbf{計算可能}（有理数演算で実行可能）
  \item 測度論的な技術的困難を回避
  \item 本質的な構造が明確化
  \item 教育的価値が向上
\end{itemize}

\subsection{構造塔理論との接続}

本実装は、5層からなる理論構築の最終段階である：

\begin{center}
\begin{tikzcd}[row sep=large]
  \text{DecidableEvents.lean} \arrow[d] \\
  \text{DecidableAlgebra.lean} \arrow[d] \\
  \text{DecidableFiltration.lean} \arrow[d] \\
  \text{ComputableStoppingTime.lean} \arrow[d] \\
  \text{\textbf{AlgorithmicMartingale.lean}}
\end{tikzcd}
\end{center}

\begin{table}[h]
\centering
\caption{構造塔と確率論の対応}
\begin{tabular}{ll}
\hline
構造塔の概念 & 確率論の概念 \\
\hline
carrier & 事象の全体 \\
Index $= \mathbb{N}$ & 時刻 \\
layer $n$ & 時刻 $n$ の可観測事象族 $\mathcal{F}_n$ \\
monotone & 情報の単調増加 $\mathcal{F}_0 \subseteq \mathcal{F}_1 \subseteq \cdots$ \\
minLayer($A$) & 事象 $A$ が初めて可観測になる時刻 \\
停止時間 $\tau$ & $\{\tau \leq t\}$ の minLayer $\leq t$ \\
\hline
\end{tabular}
\end{table}

\section{確率質量関数と期待値}

\subsection{有限標本空間}

\begin{definition}[有限標本空間]
有限標本空間 $\Omega$ は、以下を備えた型：
\begin{itemize}
  \item $\Omega$.carrier : 標本点の集合
  \item Fintype インスタンス（有限性）
  \item DecidableEq インスタンス（等号判定可能性）
\end{itemize}
\end{definition}

\subsection{確率質量関数}

\begin{definition}[確率質量関数]
有限標本空間 $\Omega$ 上の\textbf{確率質量関数} (Probability Mass Function) は、
以下を満たす関数 $p : \Omega \to \mathbb{Q}$：
\begin{enumerate}
  \item \textbf{非負性}：$\forall \omega \in \Omega,\ 0 \leq p(\omega)$
  \item \textbf{正規化}：$\sum_{\omega \in \Omega} p(\omega) = 1$
\end{enumerate}
\end{definition}

\begin{remark}[なぜ有理数か？]
有理数 $\mathbb{Q}$ を用いる理由：
\begin{itemize}
  \item \textbf{計算可能性}：exact な算術演算が可能
  \item \textbf{決定可能性}：等号判定、大小比較が decidable
  \item \textbf{形式化の容易性}：実数の完備性に依存しない
  \item \textbf{実用性}：有限標本空間では十分
\end{itemize}
実数への拡張は自然な埋め込み $\mathbb{Q} \hookrightarrow \mathbb{R}$ により可能。
\end{remark}

\subsection{期待値}

\begin{definition}[期待値]
確率質量関数 $P$ と可測関数 $X : \Omega \to \mathbb{Q}$ に対し、
$X$ の\textbf{期待値}は：
\[
  \mathbb{E}_P[X] := \sum_{\omega \in \Omega} P(\omega) \cdot X(\omega).
\]
\end{definition}

\begin{proposition}[期待値の基本性質]
期待値は以下の性質を満たす：
\begin{enumerate}
  \item \textbf{定数関数}：$X(\omega) \equiv c$ ならば $\mathbb{E}[X] = c$
  \item \textbf{線形性（加法性）}：$\mathbb{E}[X + Y] = \mathbb{E}[X] + \mathbb{E}[Y]$
  \item \textbf{線形性（斉次性）}：$\mathbb{E}[cX] = c \cdot \mathbb{E}[X]$
  \item \textbf{有限和との交換}：$\mathbb{E}\left[\sum_{i \in I} X_i\right] = \sum_{i \in I} \mathbb{E}[X_i]$ \quad ($I$ 有限)
\end{enumerate}
\end{proposition}

\begin{proof}[証明のスケッチ]
(1) 定数関数の場合：
\begin{align*}
  \mathbb{E}[c] &= \sum_{\omega} P(\omega) \cdot c \\
  &= c \sum_{\omega} P(\omega) \\
  &= c \cdot 1 = c.
\end{align*}
(2)-(4) は和の性質と Fubini の定理（有限版）から従う。
Lean 4 実装では、これらを \texttt{expected\_const}, \texttt{expected\_add},
\texttt{expected\_mul\_const}, \texttt{expected\_finset\_sum} として証明している。
\end{proof}

\subsection{指示関数と確率}

\begin{definition}[指示関数]
事象 $A \subseteq \Omega$ の\textbf{指示関数} $\mathbf{1}_A : \Omega \to \mathbb{Q}$ は：
\[
  \mathbf{1}_A(\omega) := \begin{cases}
    1 & \text{if } \omega \in A, \\
    0 & \text{if } \omega \notin A.
  \end{cases}
\]
\end{definition}

\begin{proposition}[指示関数の期待値]
$\mathbb{E}[\mathbf{1}_A] = P(A)$。
\end{proposition}

\begin{proof}
\begin{align*}
  \mathbb{E}[\mathbf{1}_A] &= \sum_{\omega \in \Omega} P(\omega) \cdot \mathbf{1}_A(\omega) \\
  &= \sum_{\omega \in A} P(\omega) \cdot 1 + \sum_{\omega \notin A} P(\omega) \cdot 0 \\
  &= \sum_{\omega \in A} P(\omega) = P(A).
\end{align*}
\end{proof}

\section{離散時間過程とマルチンゲール}

\subsection{単純過程}

\begin{definition}[単純過程]
離散時間の\textbf{単純過程} (simple process) は関数：
\[
  M : \mathbb{N} \times \Omega \to \mathbb{Q},
\]
すなわち $M = (M_n)_{n \geq 0}$ で、各 $M_n : \Omega \to \mathbb{Q}$ はランダム変数。
\end{definition}

\begin{example}[定数過程]
$M_n(\omega) \equiv c$ （すべての $n, \omega$ で一定）なる過程を\textbf{定数過程}という。
\end{example}

\subsection{適合性}

本実装では、適合性（adaptedness）を簡略化している：

\begin{definition}[適合性（簡略版）]
過程 $M$ がフィルトレーション $(\mathcal{F}_n)$ に\textbf{適合している}とは、
（現バージョンでは常に真）。
\end{definition}

\begin{remark}
完全な定義では「$M_n$ が $\mathcal{F}_n$-可測」という条件が必要。
将来の拡張として、値域側の有限代数を用いた本格的な可測性を導入予定。
\end{remark}

\subsection{マルチンゲール条件}

\begin{definition}[マルチンゲール（簡略版）]
過程 $M$ が\textbf{マルチンゲール}であるとは：
\begin{enumerate}
  \item $M$ は $(\mathcal{F}_n)$ に適合
  \item \textbf{期待値保存}：すべての $n$ で
  \[
    \mathbb{E}[M_{n+1}] = \mathbb{E}[M_n]
  \]
\end{enumerate}
\end{definition}

\begin{remark}
古典的には条件付き期待値を用いて：
\[
  \mathbb{E}[M_{n+1} \mid \mathcal{F}_n] = M_n \quad \text{(a.s.)}
\]
と定義される。本実装は全期待値のレベルでの条件（より弱い）を採用。
\end{remark}

\begin{definition}[強マルチンゲール条件]
より強い条件として、すべての $\mathcal{F}_n$-可測事象 $A$ に対し：
\[
  \mathbb{E}[M_{n+1} \cdot \mathbf{1}_A] = \mathbb{E}[M_n \cdot \mathbf{1}_A].
\]
これは局所的な期待値保存を意味し、OST の証明に必要。
\end{definition}

\subsection{具体例：定数過程}

\begin{theorem}[定数過程はマルチンゲール]
定数過程 $M_n(\omega) \equiv c$ はマルチンゲールである。
\end{theorem}

\begin{proof}
任意の $n$ に対し：
\[
  \mathbb{E}[M_{n+1}] = \mathbb{E}[c] = c = \mathbb{E}[M_n].
\]
\end{proof}

\begin{proposition}[マルチンゲールの期待値は一定]
$M$ がマルチンゲールならば、すべての $n, m \leq T$ (time horizon) に対し：
\[
  \mathbb{E}[M_n] = \mathbb{E}[M_m] = \mathbb{E}[M_0].
\]
\end{proposition}

\begin{proof}
隣接時刻での等式 $\mathbb{E}[M_n] = \mathbb{E}[M_{n+1}]$ を
帰納的につなぐことで得られる。
Lean 4 実装では \texttt{martingale\_expectation\_const} として証明。
\end{proof}

\section{停止時間と停止過程}

\subsection{停止時間}

\begin{definition}[停止時間]
$\tau : \Omega \to \mathbb{N}$ が\textbf{停止時間}であるとは、
すべての $t$ に対し事象 $\{\tau \leq t\}$ が $\mathcal{F}_t$-可測。
\end{definition}

\begin{remark}[直感的解釈]
「時刻 $t$ において、$\tau \leq t$ かどうかを判定できる」＝
停止するかどうかは現在までの情報のみで決定可能（未来を見ない）。
\end{remark}

\begin{example}[停止時間の例]
\begin{itemize}
  \item \textbf{定数停止時間}：$\tau(\omega) \equiv k$ （すべての $\omega$ で同じ時刻）
  \item \textbf{初到達時刻}：$\tau = \inf\{n : M_n \in A\}$ （集合 $A$ への初到達）
  \item \textbf{有界停止時間}：$\tau(\omega) \leq N$ for all $\omega$
\end{itemize}
\end{example}

\subsection{停止過程}

\begin{definition}[停止過程]
停止時間 $\tau$ による\textbf{停止過程} (stopped process) は：
\[
  M^\tau_n(\omega) := M_{\min(n, \tau(\omega))}(\omega).
\]
\end{definition}

\begin{remark}[直感的意味]
$M^\tau$ は「時刻 $\tau$ で過程を凍結したもの」：
\[
  M^\tau_n = \begin{cases}
    M_n & \text{if } n < \tau, \\
    M_\tau & \text{if } n \geq \tau.
  \end{cases}
\]
\end{remark}

\begin{figure}[h]
\centering
\begin{tikzpicture}[scale=0.8]
  % 時間軸
  \draw[->] (0,0) -- (10,0) node[right] {$n$};
  \draw[->] (0,0) -- (0,4) node[above] {$M_n$};

  % マルチンゲールのパス
  \draw[thick,blue] plot[smooth] coordinates {(0,1) (1,1.5) (2,2) (3,2.3) (4,2.5) (5,2.2) (6,2.4) (7,2.6) (8,2.3) (9,2.5)};

  % 停止時刻
  \draw[red,thick,dashed] (5,0) -- (5,2.2);
  \node[below] at (5,0) {$\tau$};

  % 停止過程
  \draw[thick,orange,->] plot[smooth] coordinates {(0,1) (1,1.5) (2,2) (3,2.3) (4,2.5) (5,2.2)};
  \draw[thick,orange,dashed] (5,2.2) -- (9,2.2);

  % ラベル
  \node[blue,above] at (7,2.6) {$M_n$};
  \node[orange,above] at (7,2.2) {$M^\tau_n$};
\end{tikzpicture}
\caption{停止過程の概念図}
\end{figure}

\subsection{事象 $\{\tau = n\}$ の分割}

OST の証明において重要な役割を果たすのが、標本空間の分割：

\begin{equation}
  \Omega = \bigsqcup_{n=0}^{N} \{\tau = n\},
\end{equation}
ここで $N$ は time horizon または $\tau$ の上界。

\begin{lemma}[期待値の分割]
有界停止時間 $\tau \leq N$ に対し：
\[
  \mathbb{E}[M_\tau] = \sum_{n=0}^{N} \mathbb{E}[M_n \cdot \mathbf{1}_{\{\tau = n\}}].
\]
\end{lemma}

\begin{proof}
点ごとに：
\[
  M_\tau(\omega) = \sum_{n=0}^{N} M_n(\omega) \cdot \mathbf{1}_{\{\tau = n\}}(\omega)
  \quad (\text{ただ1項のみ非零})
\]
であるから、期待値の線形性により従う。
Lean 4 では \texttt{expected\_atStopping\_as\_sum} として証明。
\end{proof}

\section{Optional Stopping Theorem}

\subsection{主定理}

\begin{maintheorem}[Optional Stopping Theorem]
$M$ を強マルチンゲールとし、$\tau$ を有界停止時間 $(\tau \leq T)$ とする。
このとき：
\[
  \mathbb{E}[M_\tau] = \mathbb{E}[M_0].
\]
\end{maintheorem}

\subsection{証明の構造}

証明は以下の4ステップからなる：

\begin{keyidea}[証明の方針]
\begin{enumerate}
  \item \textbf{テレスコーピング分解}：$M_\tau - M_0$ を増分の和に展開
  \item \textbf{期待値への持ち上げ}：点ごとの等式を期待値レベルに
  \item \textbf{各項の消去}：強マルチンゲール条件により各増分の寄与が 0
  \item \textbf{結論}：$\mathbb{E}[M_\tau - M_0] = 0$ より $\mathbb{E}[M_\tau] = \mathbb{E}[M_0]$
\end{enumerate}
\end{keyidea}

\subsection{Step 1: テレスコーピング分解}

\begin{lemma}[増分分解]
任意の $\omega$ に対し：
\[
  M_\tau(\omega) - M_0(\omega) =
  \sum_{n=0}^{T-1} \bigl(M_{n+1}(\omega) - M_n(\omega)\bigr) \cdot \mathbf{1}_{\{\tau > n\}}(\omega).
\]
\end{lemma}

\begin{proof}[証明のアイデア]
右辺の和は、$\tau(\omega) = k$ のとき：
\begin{align*}
  &\sum_{n=0}^{T-1} (M_{n+1} - M_n) \cdot \mathbf{1}_{\{k > n\}} \\
  &= \sum_{n=0}^{k-1} (M_{n+1} - M_n) \quad (\text{$n \geq k$ で指示関数は 0})\\
  &= (M_1 - M_0) + (M_2 - M_1) + \cdots + (M_k - M_{k-1}) \\
  &= M_k - M_0 = M_\tau - M_0. \quad \text{(telescoping)}
\end{align*}
Lean 4 実装：\texttt{pathwise\_increment\_decomp}, \texttt{stopped\_difference\_as\_sum}。
\end{proof}

\subsection{Step 2: 期待値への持ち上げ}

\begin{lemma}[期待値への適用]
\begin{align*}
  \mathbb{E}[M_\tau - M_0]
  &= \mathbb{E}\left[\sum_{n=0}^{T-1} (M_{n+1} - M_n) \cdot \mathbf{1}_{\{\tau > n\}}\right] \\
  &= \sum_{n=0}^{T-1} \mathbb{E}[(M_{n+1} - M_n) \cdot \mathbf{1}_{\{\tau > n\}}].
\end{align*}
\end{lemma}

\begin{proof}
期待値の線形性と有限和との交換により従う。
\end{proof}

\subsection{Step 3: 各項の消去（核心部分）}

\begin{lemma}[増分の期待値が 0]
強マルチンゲール条件の下、各 $n < T$ に対し：
\[
  \mathbb{E}[(M_{n+1} - M_n) \cdot \mathbf{1}_{\{\tau > n\}}] = 0.
\]
\end{lemma}

\begin{proof}[証明の概略]
$A := \{\tau > n\} = \{\tau \leq n\}^c$ とおく。
停止時間の定義より $\{\tau \leq n\} \in \mathcal{F}_n$ であるから、
その補集合 $A$ も $\mathcal{F}_n$-可測。

強マルチンゲール条件により：
\[
  \mathbb{E}[M_{n+1} \cdot \mathbf{1}_A] = \mathbb{E}[M_n \cdot \mathbf{1}_A].
\]

したがって：
\begin{align*}
  &\mathbb{E}[(M_{n+1} - M_n) \cdot \mathbf{1}_A] \\
  &= \mathbb{E}[M_{n+1} \cdot \mathbf{1}_A] - \mathbb{E}[M_n \cdot \mathbf{1}_A] \\
  &= 0.
\end{align*}

Lean 4 実装：\texttt{increment\_zero}（約120行の詳細な証明）。
\end{proof}

\subsection{Step 4: 結論}

以上より：
\begin{align*}
  \mathbb{E}[M_\tau] - \mathbb{E}[M_0]
  &= \mathbb{E}[M_\tau - M_0] \\
  &= \sum_{n=0}^{T-1} \mathbb{E}[(M_{n+1} - M_n) \cdot \mathbf{1}_{\{\tau > n\}}] \\
  &= \sum_{n=0}^{T-1} 0 = 0.
\end{align*}

ゆえに $\mathbb{E}[M_\tau] = \mathbb{E}[M_0]$。\qed

\subsection{証明のフロー図}

\begin{figure}[h]
\centering
\begin{tikzcd}[row sep=large, column sep=large]
  M_\tau - M_0 \arrow[d, "\text{テレスコーピング}"] \\
  \sum_{n=0}^{T-1} (M_{n+1} - M_n) \cdot \mathbf{1}_{\{\tau > n\}}
    \arrow[d, "\mathbb{E}[\cdot]"] \\
  \sum_{n=0}^{T-1} \mathbb{E}[(M_{n+1} - M_n) \cdot \mathbf{1}_{\{\tau > n\}}]
    \arrow[d, "\text{各項} = 0"] \\
  0 \arrow[d, Rightarrow] \\
  \mathbb{E}[M_\tau] = \mathbb{E}[M_0]
\end{tikzcd}
\caption{OST 証明の論理的流れ}
\end{figure}

\section{形式化の詳細}

\subsection{Lean 4 実装の構成}

\begin{computation}[実装の階層構造]
\begin{verbatim}
AlgorithmicMartingale.lean (約1460行)
  ├─ Part 1: ProbabilityMassFunction (94行)
  │   ├─ expected (期待値の定義)
  │   ├─ expected_const (定数関数)
  │   ├─ expected_add (線形性)
  │   └─ expected_finset_sum (有限和との交換)
  │
  ├─ Part 2: Martingale (390行)
  │   ├─ IsMartingale (簡略版)
  │   ├─ IsMartingaleStrong (完全版)
  │   └─ constProcess_isMartingale (具体例)
  │
  ├─ Part 3: Stopped Process (150行)
  │   ├─ stopped (定義)
  │   ├─ eventTauEq (事象 {τ = n})
  │   └─ expected_atStopping_as_sum (分割公式)
  │
  └─ Part 4: Optional Stopping Theorem (800行)
      ├─ pathwise_increment_decomp (テレスコーピング)
      ├─ increment_zero (核心補題)
      ├─ optionalStopping_global (主定理)
      └─ optionalStopping_theorem (最終形)
\end{verbatim}
\end{computation}

\subsection{核心補題の証明統計}

\begin{table}[h]
\centering
\caption{主要補題の複雑度}
\begin{tabular}{lcc}
\hline
補題名 & 行数 & 主な技法 \\
\hline
\texttt{expected\_finset\_sum} & 23 & Finset帰納法 \\
\texttt{pathwise\_increment\_decomp} & 76 & telescoping, 自然数帰納法 \\
\texttt{increment\_zero} & 129 & 補集合の可観測性, 期待値の線形性 \\
\texttt{stopped\_difference\_as\_sum} & 97 & filter, 場合分け \\
\texttt{optionalStopping\_global} & 137 & 上記の統合 \\
\hline
\end{tabular}
\end{table}

\subsection{計算可能性の実現}

\begin{computation}[実行可能な例]
Unit 空間（1点のみ）での例：
\begin{verbatim}
def unitSpace : FiniteSampleSpace := ...
def unitPMF : ProbabilityMassFunction unitSpace :=
  { pmf := fun _ => 1, ... }

def unitProcess : SimpleProcess unitSpace :=
  fun n _ => (n : ℚ)  -- 時刻に等しい値

#eval expected unitPMF (unitProcess 0)  -- 結果: 0
#eval expected unitPMF (unitProcess 1)  -- 結果: 1
\end{verbatim}
\end{computation}

\section{数学的意義と展望}

\subsection{本形式化の意義}

\begin{enumerate}
\item \textbf{計算可能性の保証}
  \begin{itemize}
    \item すべての演算が \texttt{\#eval} で実行可能
    \item 抽象的な測度論を回避
    \item 具体例での検証が可能
  \end{itemize}

\item \textbf{構造塔理論との統合}
  \begin{itemize}
    \item フィルトレーション $\leftrightarrow$ 構造塔の層
    \item 停止時間 $\leftrightarrow$ minLayer による特徴づけ
    \item 異なる数学分野の統一的理解
  \end{itemize}

\item \textbf{教育的価値}
  \begin{itemize}
    \item 有限性により本質が明確化
    \item 測度論的技術困難の回避
    \item 学部生レベルでの理解が可能
  \end{itemize}

\item \textbf{形式化の完全性}
  \begin{itemize}
    \item \texttt{sorry} なし（すべて証明済み）
    \item 型レベルでの安全性保証
    \item 自動検証による正しさの保証
  \end{itemize}
\end{enumerate}

\subsection{今後の拡張方向}

\begin{enumerate}
\item \textbf{条件付き期待値の完全実装}
  \begin{itemize}
    \item $\mathcal{F}_n$ の atom の明示的構成
    \item $\mathbb{E}[\cdot \mid \mathcal{F}_n]$ の計算可能な定義
    \item 局所マルチンゲール条件の完全版
  \end{itemize}

\item \textbf{他のマルチンゲール定理}
  \begin{itemize}
    \item Doob の最大不等式
    \item マルチンゲール収束定理（離散版）
    \item Azuma-Hoeffding の不等式
  \end{itemize}

\item \textbf{応用例の充実}
  \begin{itemize}
    \item 対称ランダムウォーク
    \item 賭博戦略の解析
    \item 最適停止問題
  \end{itemize}

\item \textbf{実数への拡張}
  \begin{itemize}
    \item $\mathbb{Q} \hookrightarrow \mathbb{R}$ の埋め込み
    \item 測度論的 OST への接続
    \item 連続時間への一般化の準備
  \end{itemize}
\end{enumerate}

\subsection{関連研究との比較}

\begin{table}[h]
\centering
\caption{マルチンゲール理論の形式化の比較}
\begin{tabular}{lccc}
\hline
特徴 & 本実装 & 測度論版 & 古典的教科書 \\
\hline
標本空間 & 有限 & 任意 & 通常は無限 \\
確率値 & $\mathbb{Q}$ & $\mathbb{R}$ & $\mathbb{R}$ \\
計算可能性 & \texttt{\#eval} 可能 & 不可 & N/A \\
形式化 & Lean 4 & Lean 4/Isabelle & 非形式的 \\
教育的価値 & 高 & 中 & 高 \\
一般性 & 低 & 高 & 高 \\
\hline
\end{tabular}
\end{table}

\section{結論}

本稿では、離散有限標本空間上での計算可能マルチンゲール理論と
Optional Stopping Theorem の完全な形式化を解説した。

\subsection{主な成果}

\begin{itemize}
\item 確率論の核心定理（OST）の完全な形式的証明（\texttt{sorry} なし）
\item 構造塔理論との明示的な対応関係の確立
\item すべての演算が計算可能（有理数ベース）
\item 約1460行の Lean 4 コードによる実装
\end{itemize}

\subsection{形式化の教訓}

\begin{enumerate}
\item \textbf{制約は明確化をもたらす}
  \begin{itemize}
    \item 有限性により本質的構造が明らかに
    \item 測度論的困難を回避しつつ核心を捉える
  \end{itemize}

\item \textbf{計算可能性の価値}
  \begin{itemize}
    \item 抽象論と具体例の往復が容易
    \item 検証可能な形での理解
  \end{itemize}

\item \textbf{段階的構築の重要性}
  \begin{itemize}
    \item 5層構造による理論の積み上げ
    \item 各層での概念の明確な役割
  \end{itemize}
\end{enumerate}

\subsection{最後に}

本形式化は、Bourbaki の構造主義と現代的な計算可能性理論を統合した一例である。
離散/有限という制約の中で、確率論の美しい構造が明瞭に浮かび上がることを示した。

この成果が、形式数学と教育、理論と実装の架け橋となることを期待する。

\vspace{1em}

\begin{center}
\fbox{%
\begin{minipage}{0.9\textwidth}
\textbf{生成情報}
\begin{itemize}
\item Lean 4 実装：CODEX (OpenAI) により生成
\item \LaTeX 解説文書：Claude Code により生成
\item 生成日時：\today
\item ソースコード：\texttt{AlgorithmicMartingale.lean} (1460行)
\end{itemize}
\end{minipage}%
}
\end{center}

\begin{thebibliography}{99}

\bibitem{williams1991}
Williams, D. (1991).
\textit{Probability with Martingales}.
Cambridge University Press.

\bibitem{durrett2019}
Durrett, R. (2019).
\textit{Probability: Theory and Examples} (5th ed.).
Cambridge University Press.

\bibitem{kallenberg2002}
Kallenberg, O. (2002).
\textit{Foundations of Modern Probability} (2nd ed.).
Springer.

\bibitem{bourbaki1968}
Bourbaki, N. (1968).
\textit{Éléments de mathématique: Théorie des ensembles}.
Hermann.

\bibitem{lean4}
de Moura, L., \& Ullrich, S. (2021).
\textit{The Lean 4 Theorem Prover and Programming Language}.
In CADE 28.

\bibitem{mathlib}
The mathlib Community (2020).
\textit{The Lean Mathematical Library}.
In CPP 2020.

\end{thebibliography}

\appendix

\section{補足：Lean 4 コードの主要部分}

\subsection{確率質量関数の定義}

\begin{verbatim}
structure ProbabilityMassFunction (Ω : FiniteSampleSpace) where
  pmf : Ω.carrier → ℚ
  nonneg : ∀ ω, 0 ≤ pmf ω
  sum_one : (∑ ω, pmf ω) = 1
\end{verbatim}

\subsection{期待値の定義}

\begin{verbatim}
def expected (P : ProbabilityMassFunction Ω)
    (X : Ω.carrier → ℚ) : ℚ :=
  ∑ ω, P.pmf ω * X ω
\end{verbatim}

\subsection{マルチンゲール条件}

\begin{verbatim}
structure IsMartingaleStrong {Ω : FiniteSampleSpace}
    (P : ProbabilityMassFunction Ω)
    (ℱ : DecidableFiltration Ω)
    (M : SimpleProcess Ω) : Prop where
  adapted : IsAdapted ℱ M
  fair_local :
    ∀ {n : ℕ} (hn : n ≤ ℱ.timeHorizon)
      (hn1 : n + 1 ≤ ℱ.timeHorizon)
      {A : Event Ω.carrier},
      A ∈ (ℱ.observableAt n hn).events →
      expected P (fun ω => M (n + 1) ω *
                           (if ω ∈ A then 1 else 0)) =
      expected P (fun ω => M n ω *
                           (if ω ∈ A then 1 else 0))
\end{verbatim}

\subsection{Optional Stopping Theorem}

\begin{verbatim}
theorem optionalStopping_theorem
    {Ω : FiniteSampleSpace}
    (P : ProbabilityMassFunction Ω)
    (ℱ : DecidableFiltration Ω)
    (M : SimpleProcess Ω)
    (hMart : IsMartingaleStrong P ℱ M)
    (τ : ComputableStoppingTime ℱ)
    (hBound : ∀ ω, τ.time ω ≤ ℱ.timeHorizon) :
  P.expected (M 0) =
    P.expected (fun ω => M (τ.time ω) ω) := by
  ...
\end{verbatim}

\end{document}
